\section{模型的评价与推广}

\subsection{模型优点}

本文以可追溯数据贯通四个问题。问题一保留5个真实观测年份，不通过插值扩大样本，并同时报告模型结构与饱和上限对远期预测的影响；问题二只对同一综述中口径一致的15项因素排序，将治疗退出、生理代偿和管理障碍留在机制层解释，避免混合不可比证据；问题三将排序结果转化为分画像候选方案，先作指南范围初筛，再比较效果、执行和坚持性；问题四明确区分官方人口预测、本文模型输出和政策情景假设，不把情景结果表述为已发生的政策效果。四套程序均设置独立数值复算或模型自检，关键结果可由支撑材料重现。

模型链条中的主要结论可回溯至输入数据、参数假设、计算公式和检验结果。Logistic和Gompertz参数对应增长速度、拐点和饱和边界；CRITIC权重来自指标离散与冲突信息；TOPSIS得分只表示规定对象内的相对优先程度；覆盖率、条件依从率和干预效果的乘积结构则区分了“提出措施”与“实际实现效果”。

\subsection{主要局限}

问题一只有5个全国观测点，2030年趋势较稳定，但2050年仍受曲线结构和饱和上限影响。问题二的排序衡量证据优先级而非因果效应，研究数量、一致率和证据等级也不能替代统一效应量。问题三的画像、复合评分和权重含专家规则，尤其时间受限型的首选稳定率仅为66.67\%；方案必须结合个体疾病、用药和基础能量需要量再作调整。问题四的覆盖、依从和效果属于透明情景参数，联合模拟区间反映设定范围内的不确定性，不是统计置信区间，也不能直接换算为慢病病例或经济损失。

\subsection{改进与推广}

获得新的全国营养调查后，可采用滚动更新或贝叶斯方法重新估计趋势参数，并按年龄、性别和城乡分层预测。若进一步收集个体随访数据，可用统一效应量替代问题二的证据代理指标，以潜类或聚类方法校准问题三画像，并验证方案的真实依从和减重效果。问题四可接入年度覆盖、依从、体重变化及成本数据，建立地区分层的动态评估模型。当前框架也可推广到高血压、糖尿病等需要“风险预测--因素识别--分层干预--政策情景”连续决策的慢性病管理问题。

% ==================== 七、参考文献 ====================
